
\documentclass{article}
\usepackage{colortbl}
\usepackage{makecell}
\usepackage{multirow}
\usepackage{supertabular}

\begin{document}

\newcounter{utterance}

\centering \large Interaction Transcript for game `wordle{-}crazy\_withclue', experiment `medium\_frequency\_words\_clue\_no\_critic', episode 10 with qwen3.5{-}27b.
\vspace{24pt}

{ \footnotesize  \setcounter{utterance}{1}
\setlength{\tabcolsep}{0pt}
\begin{supertabular}{c@{$\;$}|p{.15\linewidth}@{}p{.15\linewidth}p{.15\linewidth}p{.15\linewidth}p{.15\linewidth}p{.15\linewidth}}
    \# & \multicolumn{2}{c}{Player} && \multicolumn{2}{c}{Game Master} \\
    \hline

    \theutterance \stepcounter{utterance}  
    & & & \multicolumn{4}{p{0.6\linewidth}}{
        \cellcolor[rgb]{0.9,0.9,0.9}{
            \makecell[{{p{\linewidth}}}]{
                \texttt{\tiny{[P1$\langle$GM]}}
                \texttt{You are a language wizard who likes to guess words by using the given rules.} \\
\\ 
\texttt{Welcome to Wordle! You have six attempts to guess the target word, a valid English word of five lowercase letters (a{-}z). Please use the tags "explanation:" and "guess:" to provide a concise explanation for each guess.} \\
\\ 
\texttt{To help you make an informed guess, you will receive a clue for the word, such as} \\
\texttt{clue: snowy white.} \\
\\ 
\texttt{Here is an example guess based on the clue:} \\
\texttt{explanation: In the fairy tail Snow White, the girl is killed because she eats a poisoned apple. And the word apple has 5 letters.} \\
\texttt{guess: apple} \\
\\ 
\texttt{After each guess, your answer will be validated, and you will receive feedback indicating which letters are correct (yellow), which letters are correct but in the wrong position (black), and which letters are incorrect (purple). This feedback can be useful in determining which letters to include or exclude in your next guess.} \\
\\ 
\texttt{For example, the feedback for "apple" might be:} \\
\texttt{guess\_feedback: a<black> p<black> p<yellow> l<black> e<purple>} \\
\\ 
\texttt{The explanation should contain details about how both the guess\_feedback and clue work together in tandem to formulate a new guess.} \\
\\ 
\texttt{Do not generate any other text.} \\
\\ 
\texttt{Let's begin with your first guess.} \\
\\ 
\texttt{clue: tithing} \\
            }
        }
    }
     \\ \\

    \theutterance \stepcounter{utterance}  
    & \multicolumn{4}{p{0.6\linewidth}}{
        \cellcolor[rgb]{0.561,0.737,0.561}{
            \makecell[{{p{\linewidth}}}]{
                \texttt{\tiny{[P1$\rangle$GM]}}
                \texttt{explanation: The clue "tithing" refers to the practice of giving a tenth of one's income, often to a church. A common 5{-}letter word associated with this concept is "tithe," which means the act of giving a tenth.} \\
\texttt{guess: tithe} \\
            }
        }
    }
    & & \\ \\

    \theutterance \stepcounter{utterance}  
    & & & \multicolumn{2}{p{0.3\linewidth}}{
        \cellcolor[rgb]{0.9,0.9,0.9}{
            \makecell[{{p{\linewidth}}}]{
                \texttt{\tiny{[GM$|$GM]}}
                \texttt{attempts = 1} \\
\texttt{target\_word = tenth} \\
\texttt{target\_word\_clue = tithing} \\
\texttt{guess = tithe} \\
\texttt{guess\_feedback = t<yellow> i<purple> t<black> h<black> e<black>} \\
            }
        }
    }
    & & \\ \\

    \theutterance \stepcounter{utterance}  
    & & & \multicolumn{4}{p{0.6\linewidth}}{
        \cellcolor[rgb]{0.9,0.9,0.9}{
            \makecell[{{p{\linewidth}}}]{
                \texttt{\tiny{[P1$\langle$GM]}}
                \texttt{guess\_feedback: t<yellow> i<purple> t<black> h<black> e<black>} \\
\\ 
\texttt{Provide your response only in this format.} \\
\texttt{explanation: details} \\
\texttt{guess: word} \\
            }
        }
    }
     \\ \\

    \theutterance \stepcounter{utterance}  
    & \multicolumn{4}{p{0.6\linewidth}}{
        \cellcolor[rgb]{0.561,0.737,0.561}{
            \makecell[{{p{\linewidth}}}]{
                \texttt{\tiny{[P1$\rangle$GM]}}
                \texttt{explanation: The feedback indicates 't' is in the word but not at the start (yellow), while 'i', 'h', and 'e' are incorrect (purple/black). The second 't' is in the correct position (black). Since 'tithing' involves giving a tenth, and 'tithe' was close but had wrong letters, I need a word with 't' in the second position and likely related to the concept of a tenth or giving. However, looking at the feedback again: t<yellow> means the first 't' is wrong position, t<black> means the second 't' is correct. So the pattern is \_ T \_ \_ \_. The letters i, h, e are out. The clue "tithing" strongly suggests the word "tenth" (a tenth part). Let's check "tenth": t (wrong pos), e (out), n, t (correct pos), h (out). Wait, 'e' and 'h' were marked purple/black (incorrect). So "tenth" is invalid. Let's re{-}evaluate. The word must have 't' in the 3rd position. Letters i, h, e are excluded. The clue is "tithing". Maybe the word is "tithe" but the feedback says i, h, e are wrong? That contradicts the word "tithe". Let's re{-}read the feedback carefully: t<yellow> (1st t is in word, wrong spot), i<purple> (i is not in word), t<black> (2nd t is in word, correct spot), h<black> (h is in word, correct spot? No, black usually means correct letter wrong position in some variants, but the prompt says "correct but in the wrong position (black)" and "correct (yellow)"? Wait, the prompt says: "correct (yellow), correct but in the wrong position (black), incorrect (purple)". This is non{-}standard. Standard Wordle is Green=Correct, Yellow=Wrong Pos, Gray=Incorrect. The prompt defines: Yellow=Correct, Black=Correct but wrong pos, Purple=Incorrect.} \\
\texttt{Let's re{-}parse the feedback with the prompt's specific rules:} \\
\texttt{t<yellow>: 't' is in the word and in the correct position (1st letter).} \\
\texttt{i<purple>: 'i' is not in the word.} \\
\texttt{t<black>: 't' is in the word but in the wrong position (3rd letter).} \\
\texttt{h<black>: 'h' is in the word but in the wrong position (4th letter).} \\
\texttt{e<black>: 'e' is in the word but in the wrong position (5th letter).} \\
\\ 
\texttt{So the word has:} \\
\texttt{{-} 't' at position 1.} \\
\texttt{{-} 't' somewhere else (not pos 3).} \\
\texttt{{-} 'h' somewhere else (not pos 4).} \\
\texttt{{-} 'e' somewhere else (not pos 5).} \\
\texttt{{-} No 'i'.} \\
\\ 
\texttt{The word is 5 letters. Pattern: T \_ \_ \_ \_.} \\
\texttt{Letters available: t, h, e.} \\
\texttt{We need two more letters.} \\
\texttt{Clue: "tithing".} \\
\texttt{If the word is "tithe", the letters are t, i, t, h, e.} \\
\texttt{Feedback for "tithe":} \\
\texttt{1. t: Correct pos? Yes. (Yellow)} \\
\texttt{2. i: Not in word? But "tithe" has 'i'. If the target is "tithe", 'i' should be correct. The feedback says i<purple> (incorrect). So the target is NOT "tithe".} \\
\texttt{The target word is related to "tithing" but does not contain 'i'.} \\
\texttt{Maybe "tenth"? t, e, n, t, h.} \\
\texttt{Guess "tenth":} \\
\texttt{1. t: Correct pos? Yes. (Yellow)} \\
\texttt{2. e: In word? Yes. Pos 2. Target has e? If target is "tenth", e is at 2. If guess is "tenth", e is at 2. So e should be Yellow. But feedback for "tithe" said e<black> (wrong pos). This implies the target has 'e' but not at pos 5.} \\
\texttt{Let's reconsider the target.} \\
\texttt{Clue: "tithing".} \\
\texttt{Maybe the word is "tenet"? No.} \\
\texttt{Maybe "tutor"? No.} \\
\texttt{What 5 letter word means tithing or is related, has T at start, has T, H, E, no I?} \\
\texttt{"tithing" {-}> "tenth"?} \\
\texttt{If target is "tenth":} \\
\texttt{Guess "tithe":} \\
\texttt{t (1): Correct (Yellow).} \\
\texttt{i (2): Not in "tenth" (Purple). Correct.} \\
\texttt{t (3): In "tenth" at 4. So wrong pos (Black). Correct.} \\
\texttt{h (4): In "tenth" at 5. So wrong pos (Black). Correct.} \\
\texttt{e (5): In "tenth" at 2. So wrong pos (Black). Correct.} \\
\texttt{This matches the feedback perfectly!} \\
\texttt{Target is "tenth".} \\
\texttt{Pattern: T \_ \_ \_ \_.} \\
\texttt{We know:} \\
\texttt{Pos 1: T} \\
\texttt{Pos 2: E (from "tenth" logic, e was at 5 in guess, so it's in the word but not 5. In "tenth" it is at 2).} \\
\texttt{Pos 3: N (unknown yet, but likely).} \\
\texttt{Pos 4: T (from "tenth" logic, t was at 3 in guess, so it's in the word but not 3. In "tenth" it is at 4).} \\
\texttt{Pos 5: H (from "tenth" logic, h was at 4 in guess, so it's in the word but not 4. In "tenth" it is at 5).} \\
\texttt{So the word is T E N T H.} \\
\texttt{Let's verify "tenth" against the clue "tithing". Tithing is giving a tenth. "Tenth" is the word.} \\
\texttt{Next guess: tenth.} \\
\\ 
\texttt{explanation: The feedback indicates the first 't' is correct (yellow), while 'i' is incorrect (purple). The second 't', 'h', and 'e' are in the word but in the wrong positions (black). This suggests the word starts with 't', contains 't', 'h', and 'e', but not 'i'. The clue "tithing" refers to giving a tenth part. The word "tenth" fits the pattern: starts with 't', has 'e' (not at the end), 't' (not at the third position), and 'h' (not at the fourth position).} \\
\texttt{guess: tenth} \\
            }
        }
    }
    & & \\ \\

    \theutterance \stepcounter{utterance}  
    & & & \multicolumn{2}{p{0.3\linewidth}}{
        \cellcolor[rgb]{0.9,0.9,0.9}{
            \makecell[{{p{\linewidth}}}]{
                \texttt{\tiny{[GM$|$GM]}}
                \texttt{Error: The response should contain the 'guess:' keyword exactly once.} \\
            }
        }
    }
    & & \\ \\

    \theutterance \stepcounter{utterance}  
    & & & \multicolumn{2}{p{0.3\linewidth}}{
        \cellcolor[rgb]{0.9,0.9,0.9}{
            \makecell[{{p{\linewidth}}}]{
                \texttt{\tiny{[GM$|$GM]}}
                \texttt{game\_result = ABORT} \\
            }
        }
    }
    & & \\ \\

\end{supertabular}
}

\end{document}
